% === Begin Label Data ===
% ID: 798
% 难度星级: 
% 标签: 
% 备注: 
% 组卷引用次数: 0
% === End  Label Data ===

\begin{problem}{2011}{G}{浙江卷}{15}{概率}
某毕业生参加人才招聘会，分别向甲、乙、丙三个公司投递了个人简历．假定该毕业生得到甲公司面试的概率为 $\dfrac{2}{3}$，得到乙、丙两公司面试的概率均为 $p$，且三个公司是否让其面试是相互独立的．记 $X$ 为该毕业生得到面试的公司个数．若 $P(X=0)=\dfrac{1}{12}$，则随机变量 $X$ 的数学期望 $E(X)=$\underline{\hspace{4em}}.
\end{problem}

\begin{answer}
$\dfrac{5}{3}$
\end{answer}

\begin{solutions}
$P(X=0)=\dfrac{1}{3}(1-p)^2=\dfrac{1}{12}$，因此 $p=\dfrac{1}{2}$．

思路1：$P(X=1)=\dfrac{2}{3}(1-p)^2+\dfrac{1}{3}p(1-p)+\dfrac{1}{3}(1-p)p=\dfrac{1}{3}$，  
$P(X=2)=\dfrac{2}{3}p(1-p)+\dfrac{2}{3}(1-p)p+\dfrac{1}{3}p^2=\dfrac{5}{12}$，  
$P(X=3)=\dfrac{2}{3}p^2=\dfrac{1}{6}$．  

因此数学期望是 $E(X)=0\cdot P(X=0)+1\cdot P(X=1)+2\cdot P(X=2)+3\cdot P(X=3)=1\times\dfrac{1}{3}+2\times\dfrac{5}{12}+3\times\dfrac{1}{6}=\dfrac{5}{3}$．

思路2：利用“示性函数技巧”，假设随机变量 $X_i$ 服从两点分布 $B(1,p_i)$，其中 $p_1$、$p_2$、$p_3$ 分别表示得到甲、乙、丙公司的面试的概率，则 $X=X_1+X_2+X_3$，并且 $EX_i=P(X_i=1)=p_i$．利用数学期望的可加性，可知  
$$
E(X)=E(X_1)+E(X_2)+E(X_3)=p_1+p_2+p_3=\dfrac{2}{3}+\dfrac{1}{2}+\dfrac{1}{2}=\dfrac{5}{3}.
$$
\end{solutions}